\RequirePackage{pdfmanagement}
\SetKeys[document/metadata]{pdfversion=1.7, pdfstandard=A-2u, lang=en-GB}
\input{preamble/pdfA-compliance/glyphtounicode}
\documentclass[11pt,a4paper]{report}

\newcommand*{\PaperTitle}{A cubic weakly reversible mass-action system with an algebraic limit cycle}
\newcommand*{\PaperTitlePlain}{A cubic weakly reversible mass-action system with an algebraic limit cycle}
\newcommand*{\PaperShortTitle}{A cubic weakly reversible algebraic limit cycle}
\newcommand*{\PaperSubtitle}{chemistry-open-problems \textperiodcentered{} 18 October 2026 \textperiodcentered{} independently verified}
\newcommand*{\PaperCredits}{17-reaction network: Claude Opus 5.5 exploration agent (effort high), 17 Oct 2026 \textperiodcentered{} Rigorous certificate, 10-reaction network and manuscript: Claude Opus 5.5 (effort xhigh) \textperiodcentered{} Independent verification: Claude Opus 5.5 (effort xhigh), 18 Oct 2026}
\newcommand*{\PaperKeywords}{chemical reaction networks, weak reversibility, algebraic limit cycles, Hilbert's 16th problem, ball arithmetic}
\input{chemop-preamble}
\addbibresource{references.bib}

\begin{document}
\makepapertitle

\begin{paperabstract}
Craciun and Erban asked whether a cubic weakly reversible two-species mass-action system can have an algebraic limit
cycle, i.e.\ whether \(W^a(3)\ge1\). We answer positively with a weakly reversible network of ten reactions, all
complexes of molecularity at most three, and integer rate constants. The circle \((x-2)^2+(y-2)^2=1\) is invariant,
carries no equilibrium, and has characteristic exponent
\(\oint\operatorname{div}F\,dt\in[-3.68475192932676\pm3.3\times10^{-15}]\), certified by rigorous ball-arithmetic
integration (and confirmed independently by interval arithmetic and by a residue computation); it is therefore a
hyperbolic, stable, algebraic limit cycle.
\end{paperabstract}

\section*{The question}

Craciun and Erban~\autocite{CraciunErban2025} introduce chemical analogues of Hilbert's 16th problem. For
\(n\in\mathbb N\), \(W_n\) is the set of two-species mass-action systems whose reactions have at most \(n\) reactant and
at most \(n\) product molecules and whose reaction graph is weakly reversible \autocite[Definitions~2--3]{CraciunErban2025};
\(W^a(n)\) is the maximum number of algebraic limit cycles of systems in \(W_n\)
\autocite[Definition~5]{CraciunErban2025}. Their Table~2 records \(W^a(2)=0\), \(W^a(3)\ge0\) and \(W^a(4)\ge1\), and
they write (arXiv:2406.05057v2, §4 and §6):

\begin{openquote}
``The question of whether cubic weakly reversible two-species systems can give rise to algebraic limit cycles
remains open, leaving us with inequality \(W^{(a)}(3)\ge 0\) in Table 2.'' \dots\
``it remains an open question whether a cubic weakly reversible system can have an algebraic limit cycle.''
\end{openquote}

\begin{theorem}\label{thm:main}
\(W^a(3)\ge1\). Explicitly, the weakly reversible network with integer rate constants
\[
\begin{aligned}
&0\xrightarrow{434}2Y,\qquad 0\xrightarrow{448}X+Y,\qquad Y\xrightarrow{3374}0,\qquad 2Y\xrightarrow{1728}3Y,\\
&3Y\xrightarrow{194}Y,\qquad 3Y\xrightarrow{11}2X+Y,\qquad X+Y\xrightarrow{328}X+2Y,\\
&X+2Y\xrightarrow{250}3Y,\qquad 2X+Y\xrightarrow{216}3X,\qquad 3X\xrightarrow{28}X+2Y
\end{aligned}
\]
belongs to \(W_3\), and its mass-action system
\[
\begin{aligned}
\dot x&=f(x,y)=-56x^3+216x^2y-250xy^2+22y^3+448,\\
\dot y&=g(x,y)=56x^3-216x^2y+250xy^2+328xy-410y^3+1728y^2-3374y+1316
\end{aligned}
\]
has the circle \(\mathcal C=\{(x-2)^2+(y-2)^2=1\}\subset\mathbb R^2_{>0}\) as a hyperbolic, orbitally asymptotically
stable limit cycle.
\end{theorem}

\section{Proof}

\begin{proof}[of \zcref{thm:main}]
\emph{(1) The network is in \(W_3\).} Its eight complexes \(0, Y, 2Y, 3Y, X+Y, X+2Y, 2X+Y, 3X\) have molecularity at most
3, so every reaction has at most three reactant and at most three product molecules (condition (2.10) of
\autocite{CraciunErban2025} with \(n=3\)). The standing assumptions on networks (2.1) of \autocite{CraciunErban2025}
also hold: no reaction is repeated, no reaction has equal source and product, and both species change in some
reaction. Every reaction lies on one of the three directed cycles
\(0\to2Y\to3Y\to Y\to0\), \(0\to X+Y\to X+2Y\to3Y\to Y\to0\) and \(3Y\to2X+Y\to3X\to X+2Y\to3Y\), which together contain
all ten reactions: the network is weakly reversible, with a single (strongly connected) linkage class. The
equations follow from \(\dot z=\sum\kappa\,(\text{product}-\text{source})\,x^ay^b\) over the reactions.

\emph{(2) The circle is invariant.} With \(h(x,y)=(x-2)^2+(y-2)^2-1\), direct expansion gives the polynomial identity
\[
f\,h_x+g\,h_y=s\,h,\qquad s(x,y)=-112x^2+544xy-448x-820y^2+1728y-1008 .
\]
So \(\dot h=s\,h\) along solutions and \(\mathcal C=\{h=0\}\) is invariant. It lies in \([1,3]^2\), inside the open
positive quadrant.

\emph{(3) There is no equilibrium on \(\mathcal C\).} Parametrise \(\mathcal C\) by \(x=2+\cos\theta\),
\(y=2+\sin\theta\). The field is tangent to \(\mathcal C\), so \(F=(f,g)=P\,(-\sin\theta,\cos\theta)\) with
\(P=-(y-2)f+(x-2)g\), and \(\dot\theta=P\). The resultant
\(\operatorname{Res}_y(P,h)=75272x^4-474912x^3+909300x^2-295792x+56532\) (the leading coefficient of \(h\) in \(y\)
is~1) has no real root in \([1,3]\) (Sturm count), so \(P\) and \(h\) have no common real zero; as \(P(3,2)=792>0\)
and \(\mathcal C\) is connected, \(P>0\) on \(\mathcal C\). Equivalently and more explicitly, \(P\equiv
-2(97x^2-97xy-262x-38y+175)\pmod h\), and with \(t=\tan(\theta/2)\),
\(P=4(130t^4+135t^3+522t^2+329t+198)/(1+t^2)^2\), whose numerator has no real root (Sturm), while \(P(1,2)=520\) at
\(t=\infty\); the same method applied to \(P-18759/50\) gives the exact bound \(P>375.18\) on \(\mathcal C\).
Independently, a ball-arithmetic evaluation on 4096 subintervals of \([0,2\pi]\) gives \(P\ge348\) on \(\mathcal C\).
Consequently
\(\mathcal C\) is a periodic orbit, traversed counter-clockwise, of period
\(T=\int_0^{2\pi}d\theta/P(\theta)\in[0.0103320434082921\pm4.3\times10^{-17}]\).

\emph{(4) The periodic orbit is hyperbolic and stable.} For a periodic orbit \(\gamma\) of a planar \(C^1\) vector
field, the derivative of the Poincaré return map at \(\gamma\) equals \(\exp\bigl(\oint_\gamma\operatorname{div}F\,dt\bigr)\)
\autocite[§3.4, Theorem~2]{Perko2001}. Here
\[
\oint_{\mathcal C}\operatorname{div}F\,dt=\int_0^{2\pi}\frac{\operatorname{div}F(2+\cos\theta,2+\sin\theta)}{P(\theta)}\,d\theta
\in[-3.68475192932676\pm3.3\times10^{-15}],
\]
computed by rigorous ball-arithmetic integration (Arb's \texttt{acb\_calc\_integrate}~\autocite{Johansson2018}, via
python-flint), which returns a guaranteed enclosure. The integrand is a rational function of \(\cos\theta,\sin\theta\),
analytic on a neighbourhood of \([0,2\pi]\) because \(P\) has no zero there (step~(3)); its complex poles need no
special treatment, since on any complex ball where \(P\) may vanish the ball division returns an indeterminate value
and the integrator subdivides. Two further independent evaluations agree with this enclosure: a composite midpoint
rule in mpmath interval arithmetic, with the remainder \(\tfrac{h^3}{24}\,\phi''\) enclosed on each cell, gives
\([-3.6847574,-3.6847469]\); and, with \(t=\tan(\theta/2)\), the integral becomes that of a rational function over
\(\mathbb R\) with three simple poles in the upper half-plane (\(t=i\) and two roots of \(130t^4+\dots+198\)), and
\(2\pi i\) times the sum of its residues there is \(-3.684751929326756730\ldots\) The integral is negative, so the return map is a contraction at
\(\gamma\): \(\mathcal C\) is
a hyperbolic, orbitally asymptotically stable periodic orbit, in particular isolated among periodic orbits, i.e.\ a
limit cycle. It is algebraic, being a component of \(\{h=0\}\). As a consistency check, \(\oint_{\mathcal C}s\,dt\)
has the same enclosure, as it must: the coordinate \(h\) is transversal to \(\mathcal C\) and satisfies \(\dot h=sh\).
\end{proof}

\section{Remarks}

\begin{remark}[note=construction]
With \(h\) fixed, the conditions ``\(fh_x+gh_y=sh\)'' and ``the realisation uses only complexes of molecularity at most
3'' are linear in the rate constants and in the coefficients of \(s\); weak reversibility is enforced by a strictly
positive circulation on the support of the rate vector. A linear program found a 17-reaction network during
exploration; a mixed-integer program minimising the number of reactions (SCIP) then returned the 10-reaction support
of \zcref{thm:main}, on which the solution of the linear identity is unique up to a time rescaling. This yields the
integer rate constants. (Independent floating-point check by the verifier, \texttt{verification/milp\_min\_support.py}
with HiGHS: no network in \(W_3\) with at most nine reactions, all rate constants in \([1,10^5]\), \(\mathcal C\)
invariant and \(|P|\ge1\) with constant sign at 96 equally spaced points of \(\mathcal C\) exists; ten reactions are
attained for both orientations, and the optimal support is not unique. This bounded, sampled minimality is not used in
the proof.)
\end{remark}

\begin{remark}[note=a second example]
The 17-reaction network N17 (10 complexes, rational rate constants), whose reactions are listed in
\texttt{scripts/certify\_cycle.py}, also has \(\mathcal C\) as a stable hyperbolic limit cycle, with
\(\oint\operatorname{div}F\,dt\in[-0.636539668381478\pm2\times10^{-16}]\) and period \(1.91598\ldots\)
\end{remark}

\begin{remark}
Since \(W_3\subseteq M_3\) and \(W_3\subseteq W_4\), the example also re-proves \(M^a(3)\ge1\) and \(W^a(4)\ge1\) with a
cubic system. The network has eight complexes, one linkage class and rank two, hence deficiency five; it is not
covered by the deficiency-zero or deficiency-one theorems~\autocite{Feinberg2019}.
\end{remark}

\section{Computer verification}

\texttt{scripts/certify\_cycle.py} checks, for both networks, in exact rational arithmetic (sympy): membership in
\(W_3\) (molecularities, strong connectivity), the invariance identity and its cofactor, the resultant and the sign of
\(P\); and in rigorous ball arithmetic (python-flint/Arb): the positivity of \(P\) on 4096 subintervals and enclosures
of the period, of \(\oint s\,dt\) and of \(\oint\operatorname{div}F\,dt\). Runtime about 5~s.

\texttt{verification/verify\_independent.py}, written by the independent verifier without access to the solver's code,
re-derives the same facts by different methods: reaction-by-reaction weak reversibility and all assumptions of
(2.1)--(2.10) of \autocite{CraciunErban2025}; invariance by division and by substitution of the rational
parametrisation of \(\mathcal C\); positivity of \(P\) by a Sturm count on the tangent half-angle polynomial; rigorous
enclosures of \(\oint\operatorname{div}F\,dt\) and of the period with mpmath interval arithmetic (composite midpoint rule
with enclosed remainder) and a residue evaluation of both integrals; a floating-point check that the monodromy
determinant equals \(e^{\oint\operatorname{div}F\,dt}\approx0.0251034\); and the uniqueness, up to scaling, of the
rate vector on the 10-reaction support. Runtime about 3~min.

\paperreferences
\end{document}
